\begin{frame}{Lean 1/5: what entered the checker}
\begin{block}{Pinned source}
NBER Working Paper 22252, June 2017, 87 PDF pages. SHA-256 starts \texttt{441d0120}.
\end{block}
\begin{itemize}
\item AppliedModelingLib generated the complete \texttt{AR18RaceManMachine} paper folder.
\item The Lean target is Appendix B's algebra from (B9) and (B10) to the wage formula.
\item The code treats the differential values as real variables and takes (B9), (B10) as explicit hypotheses.
\end{itemize}
\vspace{.5em}
This is a checked \alert{support lemma}, not a proof that the economy satisfies those equations.
\end{frame}

\begin{frame}[fragile]{Lean 2/5: equation $\rightarrow$ statement $\rightarrow$ proof}
\vspace{-.5em}
\footnotesize
\[\underbrace{s_L d\ln W+(1-s_L)d\ln R=p}_{\text{(B9)}},\qquad
  \underbrace{d\ln W-d\ln R=q}_{\text{(B10)}}
  \quad\Longrightarrow\quad d\ln W=p+(1-s_L)q.\]
\vspace{-.5em}
\begin{lstlisting}
theorem wage_identity_from_B9_B10
    (s dW dR productivity relativePrice : Real)
    (hB9 : s*dW + (1-s)*dR = productivity)
    (hB10 : dW-dR = relativePrice) :
    dW = productivity + (1-s)*relativePrice := by
  linear_combination hB9 + (1-s)*hB10
\end{lstlisting}
\vspace{-.5em}
\footnotesize $s$ encodes labor's net-output share; $dW,dR$ encode the two log changes; $p$ is productivity; $q$ is the relative-price change. The slide spells Lean's $\mathbb R$ notation as \texttt{Real}. Lean verifies the algebra \emph{given} both premises.
\end{frame}

\begin{frame}{Lean 3/5: why the proof line works}
Write (B9) and (B10) as zero-valued residuals:
\begin{align*}
  A &= s\,dW+(1-s)\,dR-p,\\
  B &= dW-dR-q.
\end{align*}
The Lean proof uses $A+(1-s)B=0$. The $dR$ terms cancel and $s\,dW+(1-s)dW=dW$, leaving
\[dW-p-(1-s)q=0.\]
\vspace{.4em}
\texttt{linear\_combination} checks that weighted sum symbolically; it does not infer (B9) or (B10) from the task model.
\end{frame}

\begin{frame}[fragile]{Lean 4/5: the exact wage condition}
Set $q=-\Lambda_I/(\widehat\sigma+\varepsilon_L)$ for the binding-frontier case. The checked specialization gives
\[d\ln W=P_I-\frac{(1-s_L)\Lambda_I}{\widehat\sigma+\varepsilon_L}.\]
\small\textit{Excerpt below: binders and repeated arguments omitted for legibility.}
\begin{lstlisting}
theorem automation_wage_rises_iff ... :
    0 < dW <-> (1-s)*lambdaI/denominator
                 < productivity := by
  rw [automation_wage_change ...]
  constructor <;> intro h <;> linarith
\end{lstlisting}
\vspace{.4em}
\small The proposition is conditional on the encoded (B9) and (B10) identities. Its variables range over $\mathbb R$; positivity of $\widehat\sigma+\varepsilon_L$ is needed to derive (B10)'s sign from the economic primitives, not for this algebraic equivalence.
\end{frame}

\begin{frame}{Lean 5/5: results and open boundary}
\begin{block}{Machine results}
\texttt{lake build +AR18RaceManMachine}: passed.\\
\texttt{paper\_contribution.py check AR18RaceManMachine --fast}: passed.\\
Three support theorems use only \texttt{propext}, \texttt{Classical.choice}, and \texttt{Quot.sound}.
\end{block}
\vspace{.4em}
\begin{alertblock}{What remains open}
The full continuum equilibrium, its threshold comparative statics, and the derivation of (B9) and (B10) from the paper's model. No named result is claimed as fully proved. Source-to-Lean review and final closeout were not run.
\end{alertblock}
\end{frame}
